For \(i\in\{0,\ldots,M\}\), define
\[
R_i=\sum_{j=1}^i b_0(E_{j,0}),\qquad S_i=\sum_{j=1}^i b_1(E_{j,1}).
\]
The definition of \(\pi_{xy}\) is the pushforward of \(R\) under \((e_0,e_1)\mapsto(b_0(e_0),b_1(e_1))\), so the increments \((b_0(E_{i,0}),b_1(E_{i,1}))\) are independent and identically distributed with four-cell law \(\pi\).

We prove by induction that
\[
d_i(r,s)=\Pr\{R_i=r,S_i=s\},\qquad 0\le r<k_0,\quad 0\le s<k_1. \tag{1}
\]
For \(i=0\), both counts are zero, so (1) is the initialization \(d_0(0,0)=1\), with all other entries zero.  Suppose (1) holds at stage \(i\).  Independence of the next increment gives
\[
\Pr\{R_{i+1}=r,S_{i+1}=s\}
=\sum_{x,y\in\{0,1\}}\pi_{xy}\Pr\{R_i=r-x,S_i=s-y\}. \tag{2}
\]
Negative indices have probability zero.  Moreover, both counts have nonnegative increments, so a state whose first coordinate is at least \(k_0\), or whose second coordinate is at least \(k_1\), can never return to the retained rectangle.  Hence discarding such a state loses no contribution to a later retained entry.  Substituting the induction hypothesis in (2) proves (1) at stage \(i+1\).  Summing (1) at \(i=M\) yields
\[
\sum_{\substack{0\le r<k_0\\0\le s<k_1}}d_M(r,s)
=\Pr\{R_M<k_0,\ S_M<k_1\}. \tag{3}
\]

One scan of the \(|\mathscr E_0||\mathscr E_1|\) cells of \(R\) forms the four probabilities.  At every one of the \(M\) stages, each of at most \(k_0k_1\) retained states receives four constant-time contributions.  Two rolling arrays therefore give
\[
O(|\mathscr E_0||\mathscr E_1|+Mk_0k_1)
\]
operations and \(O(k_0k_1)\) memory.

If, for example, \(k_0=M+1\), then \(0\le R_M\le M<k_0\) identically.  The first event is certain, and summing the four-cell probabilities over \(x\) leaves the Bernoulli law of the second increment.  The same induction in one dimension gives \(O(Mk_1)\) recurrence operations and \(O(k_1)\) memory after the table scan.  The case \(k_1=M+1\) is symmetric.  If both ranks equal \(M+1\), both inequalities hold identically and the probability is one.

The adjacent exact alternative \(\texttt{lem:four-cell-multinomial-sum-alternative}\) applies to the same pushforward law \(\pi\).  Its feasible-count bijection proves that (3) also equals \(Q_{M,k_0,k_1}(\pi)\), including zero cells, and supplies the paired multinomial, recurrence, and hybrid resource accounts.  This does not alter the recurrence claims just proved; it records the second exact evaluator used by the frontier consumers.